\documentclass[10pt,a4paper]{article}
\usepackage[margin=1.7cm]{geometry}
\usepackage[T1]{fontenc}
\usepackage{lmodern,microtype,parskip,enumitem}
\usepackage{hyperref}
\hypersetup{colorlinks=true,urlcolor=blue,linkcolor=blue}
\setlength{\parskip}{4pt}
\setlength{\emergencystretch}{2em}
\setlist[itemize]{nosep,leftmargin=*}
\newcommand{\f}[1]{\texttt{\detokenize{#1}}}
\AtEndDocument{\ifnum\value{page}>1\PackageError{Week5}{Report exceeds one page}{Shorten the report.}\fi}
\begin{document}
\begin{center}
{\Large\bfseries Week 5: real-world data integration}
\end{center}

\textbf{Sources and scope.} We selected two complementary openFDA datasets relevant to our study-drug topic: \href{https://api.fda.gov/drug/ndc.json?search=generic_name:\%22METHYLPHENIDATE\%20HYDROCHLORIDE\%22\&limit=100}{100 methylphenidate product listings} from the NDC Directory (100 of 262 matches; updated 7 October 2026) and \href{https://api.fda.gov/drug/enforcement.json?search=product_description:methylphenidate\&limit=100}{66 methylphenidate recall records} from the Recall Enterprise System (all matches; updated 30 September 2026). Each has a unique source identifier. Product listings and recall events describe different facts, so neither dataset is a subset of the other. Both were accessed on 7 October 2026 without payment or registration under \href{https://open.fda.gov/license/}{openFDA's CC0 licence}. The selected fields are embedded in \f{sql/05_real_world_data.sql}; full raw JSON files are not included. Live API results may change.

\textbf{Schema and cleaning.} We added \f{FDA_Product} (product NDC as PK, generic name, labeler, dosage form, marketing start date) and \f{FDA_Recall} (recall number as PK, classification, initiation date, product description, recalling firm, recall reason). The four original tables remain unchanged and fictional. No verified mapping connects these FDA records to the mock prescriptions, so we do not invent links.
\begin{itemize}
\item No missing selected values or duplicate source keys were found in the 166 records. Empty strings become NULL; required fields reject missing values and PKs reject duplicate keys.
\item Dates are converted from YYYYMMDD to SQL DATE. All selected dates parsed successfully.
\item Names retain original spelling and punctuation: ``Pfizer Inc'' and ``Pfizer Inc.'' remain separate labels. We do not assume they identify different legal firms.
\item IDs use text, long descriptions use TEXT, and recall classifications have a CHECK constraint. JSON\_TABLE imports scalar fields; package and ingredient arrays are excluded. Repeating the import refreshes only the FDA tables inside a transaction.
\end{itemize}

\textbf{Normalization supplement to Week 2.} The new tables meet 1NF through scalar fields and one record per key. They meet 2NF under our key assumptions: each has a single-column key and no assumed composite alternative key. No further non-key dependencies are identified among the selected fields, so they meet 3NF under these rules. Repeated company names alone do not establish a dependency. No additional decomposition was needed; package-level or company details would require another normalization check.

\textbf{Queries and validation.} On MySQL 8.4.0, all 166 imported records matched the selected source values, allowing date conversion. Repeat imports preserved counts; duplicate product keys and invalid recall classifications were rejected. Week 3 operations and queries still worked, returning the expected mock results: 13 history rows, three frequent-drug summaries, two students without prescriptions and six latest-prescription rows. These queries remain demonstrations using fictional students.

Four commented adaptations in \f{sql/06_real_world_queries.sql} returned:
\begin{itemize}
\item Product overview: 100 listings.
\item Products by dosage form: seven groups, with counts 25, 21, 20, 12, 11, 7 and 4 (total 100).
\item Recall classifications: 53 Class II and 13 Class III records (total 66).
\item Latest recall per exact recalling-firm label: 14 rows, using ROW\_NUMBER with recall number to break date ties.
\end{itemize}
These results matched expectations for the saved selection; catalogue counts do not measure prescribing frequency.

\textbf{Limitations and future work.} The catalogue sample is incomplete and concerns US products; listing does not establish FDA approval. Recall records describe product issues, not student misuse or all adverse effects. We still lack diagnoses, motives, actual consumption, wellbeing, grades and medicines obtained elsewhere. The data provides product context for students and university support services, but cannot establish misuse or treatment appropriateness.

The Week 4 plans to \emph{test more examples} and \emph{check report accuracy} were addressed through import and query checks. \emph{Comparing changes over time} and \emph{asking students and staff for feedback} remain unfinished. Beyond the course, we plan voluntary surveys, motives and outcomes, and protected real-world student data; personal records would not be published. These steps address the societal question more directly.

\textbf{Reproduction.} In the Week 3 database, run \f{05_real_world_data.sql}, rerun \f{03_basic_operations.sql} and \f{04_advanced_queries.sql}, then run \f{06_real_world_queries.sql} (all in \f{sql/}). No additional importer or live API connection is needed.
\end{document}
